% ------------------------------------------------------------------------
% LaTeX - Preambel ******************************************************
% ------------------------------------------------------------------------
% pre-work
% ========================================================================
% % ToDo kennzeichnen
\newcommand{\workTodo}[1]{\textcolor{red}{todo: #1}}

% % F�r Datum und Zeit in Fusszeile
% % !!!Inhalt bei Fertigstellung der Arbeit l�schen
\newcommand{\workMarkDateTime}{\today}

% % Alle Namen werden im Titel und im hyperref-Paket eingetragen
% % !!! Ueberall f�r <Wert> das Entsprechende eintragen

 % <Typ> Studienarbeit, Dipolmarbeit, Studienarbeit oder Bachlor-Abschlussarbeit
\newcommand{\workTyp}{Technical drawings\xspace}

 % <Titel> der Arbeit
\newcommand{\workTitel}{Masterquad 2015}

 % <Studiengang> z.B. Kommunikationstechnik
\newcommand{\workStudiengang}{ASM-SB \xspace}

% <Semester> mit Jahr z.B. Sommersemester 2008  
\newcommand{\workSemester}{ASM2\xspace}

% <Name> des Studenten
\newcommand{\workNameStudentOne}{Oliver Breuning\xspace}
\newcommand{\workNameStudentTwo}{Martin Brodbeck\xspace}
\newcommand{\workNameStudentThree}{J\"urgen Schmidt\xspace}
\newcommand{\workNameStudentFour}{Phillip Woditsch\xspace}

% <Pruefer> Name des pr�fenden (betreuenden) Professor an der Hochschule
\newcommand{\workPruefer}{Prof. Dr. J\"org Friedrich\xspace} 


% %%% Nur bei Abschluss-Arbeiten

% <Datum> der Abgabe der Arbeit (Eidesstatliche Erkl�rung)
\newcommand{\workDatum}{\today\xspace}

% <Zweitpr�fer>
\newcommand{\workZweitPruefer}{\workTodo{<Zweitpr�fer>}\xspace}

% <Zeitraum>
\newcommand{\workZeitraum}{Summersemester 2015\xspace}


% %%% Nur bei Industrie-Arbeiten:

% <Firma>
\newcommand{\workFirma}{\workTodo{<Firma>}\xspace}

% <Betreuer in der Firma>
\newcommand{\workBetreuer}{\workTodo{<Betreuer in der Firma>}\xspace}

% Firmenlogo Name hier anpassen, Gr��e (wenn m�glich) nicht �ndern
%\newcommand{\workFirmenLogo}{\includegraphics[width=5cm]{fig/aa-titel/Bosch_4C_S}} 
